\subsection{GROUP GERM DEFINED BY A COMPLETE NORMED LIE ALGEBRA}

Let \(H = \sum_{r,s \geqslant 0} H_{r,s} \in \hat{L}(U, V)\) be the Hausdorff series (§ 6, no. 4, Definition 1). We shall show that the corresponding formal power series

\[
\tilde {\mathrm{H}} = \sum_ {r, s \geqslant 0} \tilde {\mathrm{H}} _ {r, s} \in \hat {\mathrm{P}} (\mathfrak {g} \times \mathfrak {g}, \mathfrak {g})\tag{3}
\]

is convergent (Differentiable and Analytic Manifolds, R, 3.1.1).

We introduce the following formal power series \(\eta\in\mathbf{Q}[[U,V]]\)

\begin{enumerate}
\def\labelenumi{(\arabic{enumi})}
\setcounter{enumi}{3}
\tightlist
\item
\item
\end{enumerate}

\[
\begin{array}{l}\eta (\mathrm{U},\mathrm{V}) = -\log (2 - \exp (\mathrm{U} + \mathrm{V}))\\ \qquad = \sum_{m\geqslant 1}\frac{1}{m} (\exp (\mathrm{U} + \mathrm{V}) - 1)^{m}\\ \qquad = \sum_{m\geqslant 1}\frac{1}{m}\sum_{\substack{r_{1},\ldots ,r_{m}\\ s_{1},\ldots ,s_{m}\\ r_{i} + s_{i}\geqslant 1}}\frac{\mathrm{U}^{r_{1}}}{r_{1}!}\frac{\mathrm{V}^{s_{1}}}{s_{1}!}\frac{\mathrm{U}^{r_{2}}}{r_{2}!}\dots \frac{\mathrm{V}^{s_{m}}}{s_{m}!}. \end{array}\tag{6}
\]

Hence

\[
\eta (\mathrm{U}, \mathrm{V}) = \sum_ {r, s \geqslant 0} \eta_ {r, s} \mathrm{U} ^ {r} \mathrm{V} ^ {s},\tag{7}
\]

where

\[
\eta_{r,s} = \sum_{m\geqslant 1}\frac{1}{m}\sum_{\substack{r_{1} + \dots +r_{m} = r\\ s_{1} + \dots +s_{m} = s\\ r_{i} + s_{i}\geqslant 1}}\frac{1}{r_{1}!\ldots r_{m}!s_{1}!\ldots s_{m}!}.\tag{8}
\]

Now let u and v be two positive real numbers such that \(u + v < \log 2\) ; then \(0 \leqslant \exp(u + v) - 1 < 1\) ; the series derived from (5) and (6) by substituting u for U and v for V are convergent and the above calculations imply that

\[
\sum_ {r, s \geqslant 0} \eta_ {r, s} u ^ {r} v ^ {s} = - \log (2 - \exp (u + v)) <   + \infty .\tag{9}
\]

Let \(r, s \geqslant 0\) and let \(\| \tilde{\mathrm{H}}_{r,s} \|\) be the norm of the continuous-polynomial \(\tilde{\mathrm{H}}_{r,s}\) (Differentiable and Analytic Manifolds, R, Appendix, no. 2).

Lemma 1.

\[
\| \widetilde {\mathrm{H}} _ {r, s} \| \leqslant \mathrm{M} ^ {r + s - 1} \eta_ {r, s}.
\]

Let \(r_i, s_i\) be in \(\mathbf{N}\) for \(1 \leqslant i \leqslant m\) , with \(s_m = 1\) ; we write \(r = \sum_{i} r_i, s = \sum_{i} s_i\) and consider the following element of \(L(\{U, V\})\) :

\[
Z = \left(\left(\sum_ {i = 1} ^ {m - 1} (\operatorname{ad} U) ^ {r _ {i}} (\operatorname{ad} V) ^ {s _ {i}}\right) (\operatorname{ad} U) ^ {r _ {m}}\right) (V).
\]

Then \(\tilde{Z} = f \circ p\) , where \(f\) is the following \((r + s)\) -linear mapping of \(\mathfrak{g}^{r + s}\) into \(\mathfrak{g}\) :

\[
\begin{array}{l} \left(x _ {1}, \dots , x _ {r}, y _ {1}, \dots , y _ {s}\right) \mapsto \\ \left(\operatorname{ad} \left(x _ {1}\right) \circ \dots \circ \operatorname{ad} \left(x _ {r _ {1}}\right) \circ \operatorname{ad} \left(y _ {1}\right) \circ \dots \circ \operatorname{ad} \left(y _ {s _ {1}}\right) \circ \operatorname{ad} \left(x _ {r _ {1} + 1}\right) \circ \dots \circ \operatorname{ad} \left(x _ {r}\right)\right) \left(y _ {s}\right) \end{array}
\]

and where \(p\) is the following mapping of \(\mathfrak{g}^2\) into \(\mathfrak{g}^{r+s}\) :

\[
(x, y) \mapsto \underbrace {(x , \dots , x} _ {r}, \underbrace {y , \dots , y)} _ {s};
\]

hence \(\| \tilde{Z}\| \leqslant \| f\| \leqslant M^{r + s - 1}\) (Differentiable and Analytic Manifolds, R, Appendix). Applying these inequalities to the various terms on the right hand side of formula (9) of §6, no. 4, we obtain:

\[
\| (\mathrm{H} _ {r, s} ^ {\prime}) ^ {\sim} \| \leqslant \frac {\mathrm{M} ^ {r + s - 1}}{r + s} \sum_ {m \geqslant 1} \frac {1}{m} \sum_ {\substack {r _ {1} + \dots + r _ {m} = r \\ s _ {1} + \dots + s _ {m - 1} = s - 1 \\ r _ {1} + s _ {1} \geqslant 1, \dots , r _ {m - 1} + s _ {m - 1} \geqslant 1}} \frac {1}{r _ {1} ! \dots r _ {m} ! s _ {1} ! \dots s _ {m - 1} !}.\tag{10}
\]

A similar argument gives

\[
\begin{array}{l}\| (\mathbf{H}_{r,s}^{\prime \prime})\sim \| \\ \leqslant \frac{\mathbf{M}^{r + s - 1}}{r + s}\sum_{m\geqslant 1}\frac{1}{m}\sum_{\substack{r_{1} + \dots +r_{m - 1} = r - 1\\ s_{1} + \dots +s_{m - 1} = s\\ r_{1} + s_{1}\geqslant 1,\ldots ,r_{m - 1} + s_{m - 1}\geqslant 1}}\frac{1}{r_{1}!\ldots r_{m - 1}!s_{1}!\ldots s_{m - 1}!} \end{array}\tag{11}
\]

whence, by (8)

\[
\left\| \tilde {\mathrm{H}} _ {r, s} \right\| \leqslant   \eta_ {r, s} \frac {\mathrm{M} ^ {r + s - 1}}{r + s} \leqslant \eta_ {r, s} \mathrm{M} ^ {r + s - 1},
\]

which proves the lemma.

PROPOSITION 1. The formal power series \(\tilde{\mathbf{H}}\) is a convergent series (Differentiable and Analytic Manifolds, R, 3.1.1); its domain of absolute convergence (Differentiable and Analytic Manifolds, R, 3.1.4) contains the open set

\[
\Omega = \left\{(x, y) \in \mathfrak {g} \times \mathfrak {g} | \| x \| + \| y \| <   \frac {1}{M} \log 2 \right\}.
\]

Let \(u, v\) be two real numbers \(> 0\) such that \(u + v < \frac{1}{M} \log 2\) ; then (Lemma 1)

\[
\begin{array}{r l} \mathrm{M} \sum_ {r, s \geqslant 0} \| \tilde {\mathrm{H}} _ {r, s} \| u ^ {r} v ^ {s} \\ & \leqslant \sum_ {r, s \geqslant 0} \eta_ {r, s} \mathrm{M} ^ {r + s} u ^ {r} v ^ {s} = - \log (2 - \exp \mathrm{M} (u + v)) <   + \infty \end{array} \tag {12}
\]

by (9).

Let \(h: \Omega \to \mathfrak{g}\) denote the analytic function (Differentiable and Analytic Manifolds, R, 3.2.9) defined by \(\tilde{\mathrm{H}}\) , that is by the formula

\[
h (x, y) = \sum_ {r, s \geqslant 0} \tilde {\mathrm{H}} _ {r, s} (x, y) = \sum_ {r, s \geqslant 0} \mathrm{H} _ {r, s} (x, y) \quad \text { for } (x, y) \in \Omega .\tag{13}
\]

This function is called the Hausdorff function of g relative to M (or simply the Hausdorff function of g if no confusion can arise). Note that \(\mathrm{H}_{r,s}(\mathrm{U}, -\mathrm{U}) = 0\) if \(r + s \geqslant 2\) and hence

\[
h (x, - x) = 0 \quad \text { for } \| x \| <   \frac {1}{2 M} \log 2.\tag{14}
\]

Similarly

\[
h (0, x) = h (x, 0) = x \quad \text { for } \| x \| <   \frac {1}{\mathrm{M}} \log 2.\tag{15}
\]

PROPOSITION 2. Let

\[
\Omega^ {\prime} = \left\{(x, y, z) \in \mathfrak {g} \times \mathfrak {g} \times \mathfrak {g} | \| x \| + \| y \| + \| z \| <   \frac {1}{M} \log \frac {3}{2} \right\}.
\]

If \((x,y,z)\in \Omega^{\prime}\) , then

\[
(x, y) \in \Omega , \quad (h (x, y), z) \in \Omega , \quad (y, z) \in \Omega , \quad (x, h (y, z)) \in \Omega\tag{16}
\]

and

\[
h (h (x, y), z) = h (x, h (y, z)).\tag{17}
\]

Let \((x,y,z)\in \Omega^{\prime}\) ; clearly \((x,y)\in \Omega\) and \((y,z)\in \Omega\) . Moreover:

\[
\| h (x, y) \| \leqslant \sum_ {r, s} \| \tilde {\mathrm{H}} _ {r, s} \| \| x \| ^ {r} \| y \| ^ {s},
\]

and hence by (13)

\[
\| h (x, y) \| \leqslant - \frac {1}{M} \log (2 - \exp M (\| x \| + \| y \|)).
\]

Now \(\mathrm{M}(\|x\| + \|y\|) < \log \frac{3}{2} - \mathrm{M}\|z\|\) ; we write \(u = \exp(\mathrm{M}\|z\|)\) ; then \(1 \leqslant u \leqslant \frac{3}{2}\) and

\[
\begin{array}{r l} \mathbf {M} (\| h (x, y) \| + \| z \|) <   - \log (2 - \exp (\log \frac {3}{2} - \mathbf {M} \| z \|)) + \mathbf {M} \| z \| \\ = - \log \left(2 - \frac {3}{2 u}\right) + \log u = \log \frac {2 u ^ {2}}{4 u - 3} \\ = \log \left(2 + \frac {2 (u - 1) (u - 3)}{4 u - 3}\right) \leqslant \log 2. \end{array}
\]

We see similarly that \((x, h(y, z)) \in \Omega\) .

We now prove (17). In the Lie algebra \(\hat{\mathbf{L}}(\{\mathbf{U},\mathbf{V},\mathbf{W}\})\)

\[
\mathrm{H} (\mathrm{H} (\mathrm{U}, \mathrm{V}), \mathrm{W}) = \mathrm{H} (\mathrm{U}, \mathrm{H} (\mathrm{V}, \mathrm{W}))
\]

by Proposition 4 of §6, no. 5. By no. 1, formula (2), we therefore have in \(\hat{\mathbf{P}}(\mathfrak{g} \times \mathfrak{g} \times \mathfrak{g}, \mathfrak{g})\) the relation

\[
\tilde {\mathrm{H}} \circ (\tilde {\mathrm{H}} \times \mathrm{Id} _ {\mathfrak {g}}) = \tilde {\mathrm{H}} \circ (\mathrm{Id} _ {\mathfrak {g}} \times \tilde {\mathrm{H}}).
\]

By Differentiable and Analytic Manifolds, R, 3.1.9, there exists a number \(\varepsilon > 0\) such that formula (17) is true when \(\|x\|\) , \(\|y\|\) and \(\|z\|\) are \(\leqslant \varepsilon\) . But the functions \((x,y,z)\mapsto h(h(x,y),z)\) and \((x,y,z)\mapsto h(x,h(y,z))\) are analytic functions on \(\Omega'\) with values in \(\mathfrak{g}\) (Differentiable and Analytic Manifolds, R, 3.2.7). As \(\Omega'\) is connected and they coincide in a neighbourhood of 0, they are equal (Differentiable and Analytic Manifolds, R, 3.2.5).

The above results imply:

Let \(\alpha\) be a real number such that \(0 < \alpha \leqslant \frac{1}{3M} \log \frac{3}{2}\) . Let

\[
\mathrm{G} = \{x \in \mathfrak {g} \mid \| x \| <   \alpha \},
\]

\(\Theta = \{(x,y) \in G \times G \mid h(x,y) \in G\}\) and \(m: \Theta \to G\) be the restriction of \(h\) to \(\Theta\) . Then:

\begin{enumerate}
\def\labelenumi{(\arabic{enumi})}
\item
  \(\Theta\) is open in \(\mathbf{G} \times \mathbf{G}\) and \(m\) is analytic.
\item
  \(x \in G\) implies \((0, x) \in \Theta\) , \((x, 0) \in \Theta\) and \(m(0, x) = m(x, 0) = x\) .
\item
  \(x \in \mathbf{G}\) implies \(-x \in \mathbf{G}, (x, -x) \in \Theta, (-x, x) \in \Theta\) and
\end{enumerate}

\[
m (x, - x) = m (- x, x) = 0.
\]

\begin{enumerate}
\def\labelenumi{(\arabic{enumi})}
\setcounter{enumi}{3}
\tightlist
\item
  Let \(x, y, z\) be elements of \(G\) such that \((x, y) \in \Theta\) , \((m(x, y), z) \in \Theta\) , \((y, z) \in \Theta\) and \((x, m(y, z)) \in \Theta\) . Then \(m(m(x, y), z) = m(x, m(y, z))\) .
\end{enumerate}

*In other words (Chapter III, §1), if we write \(-x = \sigma(x)\) , the quadruple \((G, 0, \sigma, m)\) is a Lie group germ over \(K_*\)

